% annexes.tex - the annexes of v_cvm.tex (normative)
\annex{Indication of this Recommendation's capabilities in V.8 and V.8bis}{v8}
\section{Encoding in CM and JM}
The CM and JM signals of \rec{V.8} are sequences of octets grouped in
categories (call function, modulation modes, protocols, GSTN access, PCM
modem availability, non-standard facilities).  This Recommendation needs
one bit of the modulation modes category to say ``\recname{} available'',
and an information field of two octets to say what of it.  The bit is an
allocation that ITU-T makes in a revision of \rec{V.8}; until it is made:
\begin{enumerate}[label=\alph*)]
\item the DCE shall send the information field of clause~\ref{sec:infofield}
  in the \emph{non-standard facilities} category of CM and JM, as \rec{V.8}
  provides for it, with the country code and the provider code of
  \rec{T.35} that the implementer holds, followed by the octets of
  Table~\ref{tab:infofield} with the identifier octet \code{0x43} (`C')
  first, so that a DCE of this Recommendation recognizes the field among
  a provider's other uses of the category;
\item the non-standard facilities category is the category octet
  \code{0x0F} followed by the NS field of \rec{V.8} (its length octet
  $K + L + M + 1$, the country code, the provider code's length $L$, the
  provider code, then the information: here the identifier and the two
  octets of the field), carried in extension octets of the form
  \code{xx010xxx}, five bits each, the field's bits from its first
  octet's least significant on, in bits 1 to 3 then 7 and 8 of each
  extension octet;
\item the DCE shall also set, in the modulation modes category, the bits
  of its fallback modulations (V.22 and V.21) and, in the protocols
  category, LAPM.
\end{enumerate}
On the allocation of the bit, a DCE shall set it and shall still send the
information field, in the category ITU-T assigns to it.

\section{The information field}\label{sec:infofield}
\begin{table}[htbp]
\caption{The information field}
\label{tab:infofield}
\small
\begin{tabularx}{\textwidth}{llX}
\toprule
\rowcolor{ituhead}\thd{Octet} & \thd{Bits} & \thd{Meaning} \\
\midrule
1 & 8--6 & version of this Recommendation: 001 \\
1 & 5 & signal interface at 16000~Hz (1), or 8000~Hz (0) \\
1 & 4 & facsimile of Annex~\ref{annex:fax} supported \\
1 & 3 & the amplitude loading of clause~\ref{sec:qam} supported \\
1 & 2 & the preset dictionary for \rec{V.44} of Annex~\ref{annex:v44dict} \\
1 & 1 & reserved, 0 \\
2 & 8 & robust mode (always 1) \\
2 & 7 & narrowband mode (always 1) \\
2 & 6 & wideband mode, 200~Hz comb \\
2 & 5 & wideband mode, 400~Hz comb \\
2 & 4 & 8-DPSK \\
2 & 3 & 22.5~ms symbol (the 20~ms symbol is always supported) \\
2 & 2 & quick start of clause~\ref{sec:quick} \\
2 & 1 & parametric mode (clause~\ref{sec:parametric}) \\
\bottomrule
\end{tabularx}
\end{table}
A DCE receiving an information field with a version it does not know
shall treat the peer as having version~1 and the capabilities the field's
bits it knows say.

\section{V.8bis}
A DCE that uses the procedures of \rec{V.8bis} shall carry the same
information field in the capabilities list of its MS and CL messages, in
the non-standard capability of that Recommendation, with the same
identifier octet.

\annex{The probing signals and the INFO messages}{probe}
\section{L1, L2 and L3}
L1 is each comb the transmitting DCE supports for the band of its
signal interface (Table~\ref{tab:combs}: N200, N300 and N400 at
8000~Hz; W200, W300 and W400 at 16000~Hz; S200, S300 and S400 at
32000~Hz), in the order of their codes, all tones loaded, DQPSK, with the
20~ms symbol, for 0.6~s each (30 symbols); L2 is the 200~Hz comb with the
22.5~ms symbol for 0.5~s (22 symbols).  The
phase difference transmitted on tone $k$ in symbol $n$ is given by two
bits taken from the scrambler of clause~\ref{sec:modulation} fed with
ones, started at the first symbol of L1 and run across the tones from the
lowest and across the symbols in order, so that the receiver, knowing
the sequence, measures the phase error of every tone symbol.  L1's first
symbol is the phase reference.  The level is $-12$~dBFS over all tones.

L3 is transmitted by a DCE that indicated the parametric mode
(clause~\ref{sec:parametric}): the P20 preamble and 40 symbols with the
20~ms symbol, then the P10 preamble and 80 symbols with the 10~ms
symbol, 1.84~s in all, the P10 part beginning where the P20 part ends
(the receiver takes the P10 part's timing from the P20 part's: the P10
preamble's lags are found poorly through AMR 10.2~\kbs{} and below); each symbol's index is four bits from the same
scrambler, continued from L2, so that the receiver, knowing the lags,
counts the symbols in error.  Its level is the parametric mode's,
$-15$~dBFS.

\section{What the receiver measures}
For each tone $k$ and each of L1 and L2 the receiver computes the root
mean square phase error $\sigma_k$ of the received phase differences
against the transmitted ones, after removing the mean (the clock offset's
rotation), the mean received level of the tone, the rms change
$\sigma_{A,k}$ of its magnitude from one symbol to the next in dB (for the
amplitude loading of clause~\ref{sec:qam}; L1's amplitude is constant),
and the fraction of tone symbols in error; for each part of L3, the
number of lag symbols in error.

\section{INFO0}
\begin{table}[htbp]
\caption{INFO0}
\small
\begin{tabularx}{\textwidth}{llX}
\toprule
\rowcolor{ituhead}\thd{Octet} & \thd{Bits} & \thd{Meaning} \\
\midrule
1 & 8--1 & the information field's octet 2 (Annex~\ref{annex:v8}): the modes supported \\
2 & 8--1 & the combs this DCE will probe with, a bit each for codes 0 to 7 of Table~\ref{tab:combs} (bit 8 for code 0) \\
3 & 8 & the same for code 8 \\
3 & 7 & quick start offered (the loading of octets 5 onwards is the last call's) \\
3 & 6 & quick start accepted (in the answering DCE's INFO0) \\
3 & 5 & the amplitude rings of clause~\ref{sec:qam} allowed (the DCE's \code{+CVM} permits them) \\
3 & 4--1 & reserved, 0 \\
4 & 8--1 & the DCE's transmit level in the data phases, in 0.5~dB below $-9$~dBFS \\
5-- & & with quick start: the loading of Table~\ref{tab:loading} \\
\bottomrule
\end{tabularx}
\end{table}

\section{INFO1 and MP}
\begin{table}[htbp]
\caption{INFO1 and MP: the loading asked for}
\label{tab:loading}
\small
\begin{tabularx}{\textwidth}{llX}
\toprule
\rowcolor{ituhead}\thd{Octet} & \thd{Bits} & \thd{Meaning} \\
\midrule
1 & 8 & MP only: acknowledgement (the peer's MP was received) \\
1 & 7--6 & symbol period: 00 20~ms, 01 22.5~ms, 10 10~ms (parametric mode only) \\
1 & 5--4 & code rate: 00 1/2, 01 3/4, 10 7/8 \\
1 & 3--1 & reserved, 0 \\
2 & 8 & MP only: retrain (both DCEs return to phase~2 after the exchange, clause~\ref{sec:reneg}) \\
2 & 7--1 & reserved, 0 \\
3 & 8--1 & round-trip delay this DCE measured, in 10~ms, 255 if unknown \\
4 & 8--5 & the comb: its code of Table~\ref{tab:combs}; 0 with the parametric mode \\
4 & 4 & the parametric mode (clause~\ref{sec:parametric}) is asked for, with the symbol period of octet~1: $N$ is then 0 \\
4 & 3 & amplitude rings: the second bitmap follows the first \\
4 & 2--1 & reserved, 0 \\
5 & 8--1 & number of tones $N$ in the comb, 8 to 68; 0 with the parametric mode, or when no data mode is asked for (clause~\ref{sec:loading}~e)) \\
6-- & & $N$ two-bit fields, the lowest tone first, packed from bit 8: 00 not loaded, 01 DQPSK, 10 8-DPSK, 11 reserved; where octet~4 bit~3 is set, $N$ more two-bit fields after them (from the next octet), the tones in the same order: 00 no rings, 01 two rings, 10 four rings (clause~\ref{sec:qam}), 11 reserved \\
\bottomrule
\end{tabularx}
\end{table}
Each message is an HDLC frame of \rec{V.42} (flags, the FCS of ISO/IEC
13239, no address or control octet) carrying the octets above, transmitted
twice in the message mode of clause~\ref{sec:infomod}; a receiver takes
the first copy whose FCS checks.  INFO0 is sent 100~ms after the end of
the V.8 procedures by the calling DCE; the answering DCE's INFO0 50~ms
after the end of the calling DCE's L3 (of L2 where there is no L3), and
its INFO1 100~ms after the end of its own L3; the calling DCE's INFO1
50~ms after the end of the answering DCE's.  Each DCE sends MP 100~ms
after the last INFO1, and after it has received its peer's MP with the
acknowledgement bit set, one more MP with its own set.

\annex{The inner code and its interleaver}{rs}
\section{The code}
The inner code is a Reed-Solomon code over GF($2^8$).  The field is
generated by the primitive polynomial $p(x) = x^8 + x^4 + x^3 + x^2 + 1$,
and the code generator polynomial is
\[ g(x) = \prod_{i=0}^{2t-1} (x - \alpha^i), \]
with $\alpha$ a root of $p(x)$: the field and the generator of the
RS(255,239) code of \rec{G.975} (where $t = 8$), with $t$ as
Table~\ref{tab:rsannex} gives.  Codewords are shortened to $n = 32$
octets: the encoder computes the check octets of the $k$ information
octets preceded by $255 - n$ zero octets, which are not transmitted, and
the decoder restores them.  The first octet of a codeword is the highest
order coefficient.

\begin{table}[htbp]
\caption{The code rates}
\label{tab:rsannex}
\small
\begin{tabular}{lcccc}
\toprule
\rowcolor{ituhead}\thd{Rate} & $n$ & $k$ & $t$ & \thd{erasures corrected ($2t$)} \\
\midrule
7/8 & 32 & 28 & 2 & 4 \\
3/4 & 32 & 24 & 4 & 8 \\
1/2 & 32 & 16 & 8 & 16 \\
\bottomrule
\end{tabular}
\end{table}

\section{The interleaver}
Codewords are interleaved in blocks of $D = 4$: the 4 codewords of a block
are transmitted octet by octet across the codewords, that is, octet~0 of
codeword~0, octet~0 of codeword~1, \ldots, octet~0 of codeword~3, then
octet~1 of each, and so on, 128 octets in all.  Consecutive octets on the
line thus belong to different codewords, and a symbol erasure of up to 4
octets (16 DQPSK tones) costs each codeword one octet.  The interleaver's
delay is one block: 128 octets, which is 0.64~s at 1600~bit/s and
0.31~s at 3300~bit/s.  In the parametric mode $D = 1$: the codewords are
sent one after the other (a block of 128 octets would take 5~s at
200~bit/s), and an erased frame costs the codeword it falls in one or
two octets, which every code rate corrects.  The receiver's
de-interleaver is the inverse.  The
block boundary is marked by the first symbol of phase~4 and by the MP
exchange of a renegotiation; the receiver recovers it, if lost, from the
syndromes of the codewords (a block boundary error leaves every codeword
uncorrectable).

\section{Erasure marking}
The receiver should give the decoder, per octet, an erasure flag set when
the tone symbols that carried the octet's bits were marked unreliable
(clause~\ref{sec:mapping}) or lay in a frame detected as erased
(clause~\ref{sec:erasure}).  A decoder given $e$ erasures and $v$ errors
corrects the codeword when $2v + e \le 2t$.

\annex{Facsimile}{fax}
\section{Principle}
A DCE that indicates facsimile in its information field
(Annex~\ref{annex:v8}) terminates the \rec{T.30} session of the facsimile
terminal on its side -- as a Group~3 facsimile machine does on the
PSTN, with the V.21, V.27ter, V.29 and V.17 modulations towards the
terminal, or as a T.38 gateway does towards an IP network -- and carries
the session to its peer as the Internet facsimile protocol (IFP) packets
of \rec{T.38} in the data link of this Recommendation.  The peer DCE
reproduces the session towards its own terminal.  The procedure is that
of a \rec{T.38} gateway pair with the data link of this Recommendation in
place of the IP network, and of the facsimile relay of \rec{V.150.1} in its
aim.

\section{Carriage of the IFP packets}
Each IFP packet shall be sent in the information field of one LAPM I
frame, preceded by one octet of type (\code{0x01}: T.38 IFP packet, other
values reserved) and two octets of length (most significant first).  The
UDPTL transport of \rec{T.38} and its redundancy are not used: LAPM
delivers the packets in order and without loss.  The T.38 options in use
shall be version~0, no error correction and no fill bit removal, unless
both DCEs indicated otherwise in a T.38 options exchange carried the same
way with type \code{0x02}.

\section{Rates}
The slowest image modulation of \rec{T.30} is V.27ter at 2400~bit/s.  Where
the coded rate of the data link is below the modulation rate the
terminals negotiate, the DCE shall accept only the error correction mode
of \rec{T.30} Annex~A from its terminal (it shall clear the non-ECM
capability in the DIS it relays) and shall pace the sending terminal with
the receive-not-ready procedure of that annex, as T.38 gateways do.  The
T.30 session's timers (T1, T2, T4) are met by the DCE towards its
terminal; the DCE may spoof the terminal with preamble flags while the
peer's answer is in transit, as \rec{T.38} allows.

\annex{A preset dictionary for V.44}{v44dict}
\section{Use}
Where both DCEs indicate it in the information field
(Table~\ref{tab:infofield}, octet~1 bit~2) and the XID exchange of
\rec{V.42} agrees on \rec{V.44} compression, each DCE shall, before the
first data in either direction, prime its V.44 encoder and its V.44
decoder with the text of clause~\ref{sec:v44text}, as follows:
\begin{enumerate}[label=\alph*)]
\item with the parameters the XID exchange agreed (N2, N7 and N8 of
  \rec{V.44}), the text's characters are given, in order, to a V.44
  encoder in compressed mode, from the start of the text, for as long as
  \begin{itemize}
  \item the next free codeword, plus twice the number of characters
    taken but not yet encoded, plus~6, is below N2, and
  \item the characters taken are fewer than $\lfloor N8/2 \rfloor - 1$;
  \end{itemize}
  then the encoder's FLUSH ends them. When the parameters let the whole
  text through, all of it is taken (3505 characters);
\item the dictionary and the history that a V.44 decoder holds once it
  has decoded that encoder's output are the primed state, for the
  encoder and the decoder of both directions. A DCE need not run the
  encoder: any procedure that reaches the same state will do;
\item the data then follows, encoded and decoded from the primed state
  in compressed mode, the codeword size, the ordinal size, the next free
  codeword and the history position those of the primed state;
\item a REINIT control code, sent or received, returns the dictionary
  and the history to the primed state (and the characters the encoder has
  not encoded yet go into the history after the text), in place of the
  empty ones of \rec{V.44}; so does the decoder's ECM after transparent
  mode;
\item a new XID exchange that agrees on V.44 primes them anew.
\end{enumerate}
Where either DCE does not indicate it, or the XID exchange agrees on
\rec{V.42bis} or no compression, \rec{V.44} and \rec{V.42bis} run as those
Recommendations specify.

\itunote{V.44 starts with an empty dictionary and learns from the data.
At the rates of this Recommendation's parametric mode and narrowband
combs (200 to 3500~bit/s) a call's data is often a few short messages,
over before the dictionary has learnt much. Primed, V.44 halves such
messages: six typical ones (an HTTP request and response, a mail, JSON,
NMEA sentences, a short letter, 1264 octets) come to 958 octets with V.44
alone, 485 primed. The text costs nothing on the line and about 3.5~kB of
history at each end.}

\section{The text}\label{sec:v44text}
The text is the lines below, each followed by CR LF (\code{0x0D 0x0A}),
in this order: 113 lines, 3505 octets, all of 7-bit characters
(\rec{T.50}), so that the ordinals stay 7 bits. A visible space
(\textvisiblespace) is a space (\code{0x20}); line~27 is empty.
{\footnotesize
\begin{verbatim*}
GET / HTTP/1.1
Host: www.example.com
User-Agent: Mozilla/5.0 (compatible)
Accept: text/html,application/xhtml+xml,application/xml;q=0.9,*/*;q=0.8
Accept-Language: en-US,en;q=0.5
Accept-Encoding: gzip, deflate
Connection: keep-alive
Cache-Control: no-cache
Content-Type: application/x-www-form-urlencoded
Content-Length: 0
Authorization: Basic 
Cookie: session=
POST /api/v1/ HTTP/1.1
Content-Type: application/json; charset=utf-8
HTTP/1.1 200 OK
HTTP/1.1 301 Moved Permanently
HTTP/1.1 304 Not Modified
HTTP/1.1 404 Not Found
HTTP/1.1 500 Internal Server Error
Date: Mon, 01 Jan 2024 00:00:00 GMT
Server: Apache
Last-Modified: Tue, 02 Feb 2024 12:00:00 GMT
Location: https://www.example.com/index.html
Set-Cookie: id=; Path=/; HttpOnly
Transfer-Encoding: chunked
Content-Type: text/html; charset=UTF-8

<!DOCTYPE html>
<html lang="en">
<head>
<meta charset="utf-8">
<meta name="viewport" content="width=device-width, initial-scale=1">
<title></title>
<link rel="stylesheet" href="style.css">
<script type="text/javascript" src=""></script>
</head>
<body>
<div class=""><span></span></div>
<p><a href="https://www.">.com</a></p>
<ul><li></li></ul>
<table><tr><td></td></tr></table>
<img src="" alt=""><br>
<form action="" method="post"><input type="text" name="" value="">
</body>
</html>
<?xml version="1.0" encoding="UTF-8"?>
{"id": 1, "name": "", "type": "", "value": 0, "status": "ok",
 "time": "2024-01-01T00:00:00Z", "data": [], "error": null,
 "enabled": true, "count": 0, "message": "", "result": false}
From: <user@example.com>
To: <user@example.com>
Cc: 
Subject: Re: 
Date: Wed, 03 Jan 2024 10:00:00 +0000
Message-ID: <@mail.example.com>
Reply-To: 
MIME-Version: 1.0
Content-Type: text/plain; charset="us-ascii"
Content-Type: multipart/mixed; boundary="
Content-Transfer-Encoding: 7bit
Content-Disposition: attachment; filename="
Received: from mail.example.com by mx.example.org with ESMTP id
EHLO example.com
250-SIZE 10240000
MAIL FROM:<user@example.com>
RCPT TO:<user@example.com>
DATA
354 Start mail input; end with <CRLF>.<CRLF>
250 OK
220 Service ready
221 Bye
QUIT
+OK POP3 server ready
USER 
PASS 
STAT
LIST
RETR 1
DELE 1
-ERR 
A001 LOGIN 
A002 SELECT INBOX
A003 FETCH 1 (FLAGS BODY[HEADER])
* OK [CAPABILITY IMAP4rev1] ready
$GPGGA,000000.00,0000.0000,N,00000.0000,E,1,08,1.0,0.0,M,0.0,M,,*00
$GPRMC,000000.00,A,0000.0000,N,00000.0000,E,0.0,0.0,010124,,,A*00
$GPGSV,3,1,12,
$GPVTG,0.0,T,,M,0.0,N,0.0,K,A*00
login: 
Password: 
Last login: 
Welcome to 
Connection closed by foreign host.
ERROR
OK
0123456789 10 100 1000 2024 00:00:00
https://www.example.com/index.html?id=&page=1&lang=en
http://www.example.org/ .net/ .gov/ .edu/
The 
the and of to in is that for it with as was on be at by this are from
 have or not but which you an they his her were had has their been
 one all there more if will can would what so out about up who when
 also into its other new some time only these may than then them two
 first any my like over such our well even most made after should
 could because between through during before under while where your
 please thank you regards information message number service report
 system user name address phone email account date error status
 data file page text total price order today tomorrow yesterday
ing tion ment ness able ally ed. es, ly s. 
Dear Sir or Madam,
Thank you for your message. Best regards,
-----BEGIN-----
-----END-----
\end{verbatim*}
}
